%=================================================================
% Observation Quotients, Orbit Predicates, and the Derivability of
% Control Limits -- Version 11
%=================================================================
\documentclass[11pt]{article}

\usepackage[T1]{fontenc}
\usepackage[utf8]{inputenc}
\usepackage{lmodern}
\usepackage[margin=1in]{geometry}
\usepackage{microtype}
\usepackage{mathrsfs}
\usepackage{amsmath,amssymb,amsthm,mathtools}
\usepackage{booktabs,array,longtable}
\usepackage{enumitem}
\usepackage[hidelinks]{hyperref}
\usepackage[nameinlink,capitalise]{cleveref}

\setlength{\parskip}{0.42em}
\setlength{\parindent}{0pt}
\setlist{itemsep=0.25em,topsep=0.35em}

%================== macros ==================
\newcommand{\Bbits}{\ensuremath{\{0,1\}}}
\newcommand{\Xspace}{\ensuremath{X}}
\newcommand{\Endo}{\ensuremath{T}}
\newcommand{\Obs}{\ensuremath{O}}
\newcommand{\Tail}{\ensuremath{\Psi}}
\newcommand{\Pred}{\ensuremath{\phi}}
\newcommand{\Diag}{\ensuremath{\mathsf{Diag}}}
\newcommand{\DOne}{\ensuremath{\mathsf{D1}}}
\newcommand{\DTwo}{\ensuremath{\mathsf{D2}}}
\newcommand{\DThree}{\ensuremath{\mathsf{D3}}}
\newcommand{\RCA}{\ensuremath{\mathsf{RCA}_0}}
\newcommand{\ACA}{\ensuremath{\mathsf{ACA}_0}}
\newcommand{\UGapTail}{\ensuremath{\mathsf{UGapTail}}}
\newcommand{\Mspace}{\ensuremath{\mathcal M}}
\newcommand{\Lib}{\ensuremath{\mathcal L}}
\newcommand{\Arch}{\ensuremath{\mathcal A}}
\newcommand{\Inv}{\ensuremath{\mathcal I}}
\newcommand{\FNovel}{\ensuremath{\mathsf{FNovel}}}
\newcommand{\Base}{\ensuremath{\mathfrak B}}
\newcommand{\Task}{\ensuremath{\mathsf Q}}
\newcommand{\im}{\ensuremath{\operatorname{im}}}
\newcommand{\ran}{\ensuremath{\operatorname{ran}}}
\newcommand{\Dec}{\ensuremath{\operatorname{Dec}}}
\newcommand{\Proofs}{\ensuremath{\mathsf{Prf}}}
\newcommand{\MInf}{\ensuremath{\mathcal M_\infty}}
\newcommand{\MEnc}{\ensuremath{\mathsf{Enc}}}
\newcommand{\MDec}{\ensuremath{\mathsf{Dec}}}
\newcommand{\AuditQuot}{\ensuremath{q_{\mathcal A}}}
\newcommand{\AuditEq}{\ensuremath{\sim_{\mathcal A}}}
\newcommand{\Rthree}{\ensuremath{\mathbb R^3}}
\newcommand{\Ledger}{\ensuremath{\mathscr L}}
\newcommand{\Debt}{\ensuremath{B}}
\newcommand{\Demand}{\ensuremath{W}}
\newcommand{\AuditCap}{\ensuremath{K}}


\theoremstyle{plain}
\newtheorem{theorem}{Theorem}[section]
\newtheorem{proposition}[theorem]{Proposition}
\newtheorem{lemma}[theorem]{Lemma}
\newtheorem{corollary}[theorem]{Corollary}
\theoremstyle{definition}
\newtheorem{definition}[theorem]{Definition}
\newtheorem{assumption}[theorem]{Assumption}
\newtheorem{construction}[theorem]{Construction}
\theoremstyle{remark}
\newtheorem{remark}[theorem]{Remark}
\newtheorem{warning}[theorem]{Warning}

\hypersetup{
 pdftitle={Observation Quotients, Orbit Predicates, and the Derivability of Control Limits: Version 11},
 pdfauthor={Andy E. Williams}
}

\title{Observation Quotients, Orbit Predicates, and the\\
Derivability of Control Limits\\[3pt]
\large A Typed Second-Order Diagnostic, an Exact Uniform\\
Reverse-Mathematical Principle, Universal M-Space Proof Architecture, and\
\large a Collective Proof-Audit Capacity Criterion}
\author{Andy E. Williams\thanks{Correspondence: info@cc4ci.org}\\[3pt]
\small Caribbean Center for Collective Intelligence, St.~John's, Antigua and Barbuda\\
\small Nobeah Foundation, Nairobi, Kenya}
\date{2026 -- Version 11}

\begin{document}
\maketitle

\begin{abstract}
For a declared observation map, a Boolean orbit target is decidable from the
available information exactly when it is constant on the observation fibers.  We
prove this criterion and construct a strictly refining observation-history tower
in which every finite monitoring predicate is decidable at its corresponding
finite stage while a distinct shift-invariant orbit target is not decidable even
from the full observation history.  The construction yields a finite paired
certificate of this non-descent.

The resulting contribution is a typed diagnostic acting on a declared binding of
an observation channel, an orbit target, and an evidence object.  A descent
certificate yields status $\DOne$; a finite equal-observation, unequal-target
pair yields $\DTwo$; and an incomplete record yields $\DThree$.  We prove that
this diagnostic does not factor through a component projection that preserves
available ingredients while forgetting which channel--target--certificate
binding is under assessment.

We also define the class of communicable purely mathematical proof witnesses and
prove that every member has an exact representation in a finite level of an
M-space proof tower.  Relative to a declared audit family, the associated audit
quotient is the coarsest audit-faithful representation, hence maximally
compressive in the task-relative sense that a theorem can support.  Every finite
proof-dependency graph has an edge-separated realization in $\mathbb R^3$, and
three is the least universal Euclidean display dimension for that class.  Under
an explicit finite-interface audit model and a declared finite admissible
projection plan, the quotient is structurally auditable through a finite indexed
record of coherent three-dimensional views.  For graph-shaped proof content,
the record reduces to one view only when the full labelled incidence display lies
within the declared interface bound; the unconditional graph result remains a
geometric realization theorem.  Finally, a novelty-separated derivational
problem admits no reliable solution through a channel family that collapses
target-distinct states; any reliable replacement must retain an
audit-equivalent refinement.  We also formalize a time-indexed proof-audit
ledger.  Under a declared audit-cost monotonicity condition, the audit quotient
weakly minimizes audit demand among audit-faithful representations; persistent
audit demand above available capacity forces unbounded unresolved audit debt.
If the quotient itself lies within capacity, replacing a higher-demand faithful
representation by the quotient is a sufficient representation-level repair.
A separate preregistered corpus and interface protocol is the empirical bridge
for whether these formal quantities forecast real repository burden or improve
human task performance.

Separately, we formulate the uniform gap-tail principle $\UGapTail$ and prove,
over $\RCA$, that $\UGapTail\leftrightarrow\ACA$.  The quotient, temporal,
binding, proof-architecture, and reverse-mathematical results are distinct
formal claims.  Application and physical interpretations remain bridge-required.
\end{abstract}

\noindent\textbf{Keywords:} observation quotient; factorization; orbit predicate;
finite monitoring; tail decision; reverse mathematics; arithmetical
comprehension; diagnostic binding; M-space; proof representation; audit quotient;
three-dimensional visualization; novelty-separated derivation.

\section{Motivation, scope, and result type}
\label{sec:intro}

A finite sequence of observations may support every declared finite monitoring
judgment while remaining insufficient for a different, long-run orbit judgment.
The mathematical question is prior to any application-specific decision:
\begin{quote}
Given a declared observation representation and a declared orbit predicate, does
the predicate descend through that representation?
\end{quote}

This paper answers that question at four levels.  First, it gives the general
fiber criterion.  Second, it gives a temporal witness separating finite
monitoring from an independently typed orbit target.  Third, it formalizes a
diagnostic whose inputs include a representation, a target, and a certificate.
Fourth, it applies the same descent discipline to communicable mathematical
proofs, their audit quotients, and their human-facing visual realizations.
Fifth, it separates a formal proof-audit capacity criterion from the empirical
question of whether source-visible mathematical corpora exhibit such capacity
bottlenecks.  The diagnostic is second order in the representational sense: it
acts on an assessment record rather than on a raw state alone.

The paper uses ``second-order'' in one additional, explicitly separate sense:
second-order arithmetic.  The reverse-mathematical theorem concerns only the
uniform principle $\UGapTail$ stated in Section~\ref{sec:reverse}; it does not
assert that every tail problem, application, or physical system requires
$\ACA$.  The M-space theorem package is likewise sharply typed: it concerns
finitely communicable proof witnesses, audit-relative compression, and
structural audit, rather than arbitrary semantic truth, literal notation, or an
unmodeled universal psychology of mathematical understanding.

\begin{remark}[Removal test]
All theorem statements and proofs below are independent of any proposed
application.  The bridge record of Section~\ref{sec:bridge} identifies what an
application would have to supply; it contributes no premise to the formal core.
\end{remark}

\section{Quotient decidability}
\label{sec:quotient}

\begin{definition}[Decidability through a representation]
Let $q:\Xspace\to Z$ and $\Tail:\Xspace\to\Bbits$.  Write
\[
\Dec_q(\Tail)
\]
when there exists a map $\overline\Tail:\im(q)\to\Bbits$ such that
\[
\Tail=\overline\Tail\circ q.
\]
The codomain is $\im(q)$, so no surjectivity convention is required.
\end{definition}

\begin{theorem}[Fiber criterion]
\label{thm:fiber}
For $q:\Xspace\to Z$ and $\Tail:\Xspace\to\Bbits$,
\[
\Dec_q(\Tail)
\quad\Longleftrightarrow\quad
\forall x,y\in\Xspace\;\bigl(q(x)=q(y)\Longrightarrow\Tail(x)=\Tail(y)\bigr).
\]
\end{theorem}

\begin{proof}
If $\Tail=\overline\Tail\circ q$, equal $q$-values give equal target values.
Conversely, suppose $\Tail$ is constant on every $q$-fiber.  For
$z\in\im(q)$ choose $x$ with $q(x)=z$ and set
$\overline\Tail(z)=\Tail(x)$.  Fiber constancy makes this definition
independent of the chosen representative, and then
$\Tail=\overline\Tail\circ q$.
\end{proof}

\begin{definition}[Paired non-descent certificate]
A paired non-descent certificate for $(q,\Tail)$ is a pair $(x,y)\in\Xspace^2$
with
\[
q(x)=q(y)
\qquad\text{and}\qquad
\Tail(x)\neq\Tail(y).
\]
\end{definition}

\begin{corollary}[Finite separation]
A paired non-descent certificate proves $\neg\Dec_q(\Tail)$.
\end{corollary}

\begin{proof}
It violates the right-hand side of Theorem~\ref{thm:fiber}.
\end{proof}

\begin{definition}[Decision partition]
For $\Tail:\Xspace\to\Bbits$, let
\[
x\equiv_{\Tail}y\quad\Longleftrightarrow\quad\Tail(x)=\Tail(y).\]
A representation $r:\Xspace\to R$ is decision-sufficient for $\Tail$ when
$\Dec_r(\Tail)$ holds.
\end{definition}

\begin{proposition}[Coarsest decision partition]
\label{prop:coarsest}
The quotient map $x\mapsto[x]_{\equiv_\Tail}$ is decision-sufficient for
$\Tail$.  If $r$ is decision-sufficient for $\Tail$, then
\[
r(x)=r(y)\Longrightarrow\Tail(x)=\Tail(y).
\]
Thus the truth partition of $\Tail$ is the coarsest partition sufficient to
decide $\Tail$.
\end{proposition}

\begin{proof}
The map $[x]_{\equiv_\Tail}\mapsto\Tail(x)$ is well defined.  If
$\Tail=d\circ r$, then $r(x)=r(y)$ implies
$\Tail(x)=d(r(x))=d(r(y))=\Tail(y)$.
\end{proof}

\begin{remark}[Scope]
Proposition~\ref{prop:coarsest} is a partition theorem.  It is not a memory
lower bound, an orbit-storage theorem, or an empirical claim about a controller.
\end{remark}

\section{Finite monitoring and independent orbit targets}
\label{sec:temporal}

\begin{definition}[History tower]
A history tower is a triple $(\Xspace,\Endo,o)$ with
\[
\Endo:\Xspace\to\Xspace,\qquad o:\Xspace\to\Obs.
\]
For $n\in\mathbb N$, set
\[
q_n(x)=\bigl(o(x),o(\Endo x),\ldots,o(\Endo^n x)\bigr),
\qquad
q_\infty(x)=\bigl(o(\Endo^m x)\bigr)_{m\in\mathbb N}.
\]
\end{definition}

\begin{definition}[Finite monitoring and orbit target]
Let $\Pred:\Xspace\to\Bbits$.  Define
\[
P_n(x)=1\quad\Longleftrightarrow\quad
\forall k\leq n\;\Pred(\Endo^k x)=1.
\]
An orbit target is a shift-invariant map $\Tail:\Xspace\to\Bbits$ satisfying
$\Tail(\Endo x)=\Tail(x)$.  No equality between $\Tail$ and a limit of the
$P_n$ is assumed.
\end{definition}

\begin{warning}[Two predicate families]
A fixed predicate that descends through every $q_n$ necessarily descends through
$q_\infty$, because each $q_n$ is a function of $q_\infty$.  The theorem below
therefore concerns two distinct families: the finite predicates $P_n$ and an
independently typed orbit target $\Tail$.
\end{warning}

\begin{theorem}[Rich temporal non-descent]
\label{thm:rich}
There exist a history tower $(\Xspace,\Endo,o)$, a local predicate $\Pred$, and
a shift-invariant target $\Tail$ such that:
\begin{enumerate}[label=\textnormal{(\roman*)},leftmargin=2.6em]
\item $P_n$ descends through $q_n$ for every $n$;
\item $q_{n+1}$ strictly refines $q_n$ for every $n$; and
\item $\Tail$ does not descend through $q_\infty$.
\end{enumerate}
\end{theorem}

\begin{proof}
Let
\[
\Xspace=\{A,B\}\times\Bbits^{\mathbb N}\times\mathbb N,
\]
with elements $(\theta,s,i)$ and $s=(s_0,s_1,\ldots)$.  Set
\[
\Endo(\theta,s,i)=(\theta,s,i+1),\qquad
 o(\theta,s,i)=s_i,\qquad
\Pred(\theta,s,i)=s_i.
\]
Then
\[
q_n(\theta,s,i)=(s_i,\ldots,s_{i+n})
\]
and
\[
P_n(\theta,s,i)=1\quad\Longleftrightarrow\quad
s_i=\cdots=s_{i+n}=1.
\]
Thus $P_n=\overline P_n\circ q_n$, where $\overline P_n$ tests whether every
coordinate is $1$.  Strict refinement follows by choosing two streams agreeing
on indices $i,\ldots,i+n$ and differing at $i+n+1$.

Define
\[
p(A,s,i)=0,\qquad p(B,s,i)=1-\frac{1}{i+1},
\]
and
\[
\Tail(x)=1\quad\Longleftrightarrow\quad
\limsup_{m\to\infty}p(\Endo^m x)<1.
\]
This is shift invariant.  For $\mathbf 1=(1,1,\ldots)$, put
\[
x_A=(A,\mathbf 1,0),\qquad x_B=(B,\mathbf 1,0).
\]
Then $q_\infty(x_A)=q_\infty(x_B)=\mathbf 1$, whereas
$\Tail(x_A)=1$ and $\Tail(x_B)=0$.  Theorem~\ref{thm:fiber} gives
$\neg\Dec_{q_\infty}(\Tail)$.
\end{proof}

\begin{corollary}[History-derived representations cannot repair non-descent]
\label{cor:derived}
In the witness of Theorem~\ref{thm:rich}, if $r=g\circ q_\infty$ for some map
$g$, then $\neg\Dec_r(\Tail)$.
\end{corollary}

\begin{proof}
The states $x_A$ and $x_B$ have equal $r$-values and unequal target values.
\end{proof}

\begin{definition}[Deterministic finite-state controller]
For an observation alphabet $O$, a deterministic finite-state controller is a
finite set $S$, an initial state $s_0\in S$, an update map
$u:S\times O\to S$, and an output map $v:S\to\Bbits$.  On the orbit of $x$,
its state sequence is
\[
s_{m+1}=u\bigl(s_m,o(T^m x)\bigr).
\]
It decides an orbit target $\Psi$ in the eventual-output sense when, for every
$x$, there is $N$ such that $v(s_m)=\Psi(x)$ for all $m\geq N$.
\end{definition}

\begin{corollary}[Finite-state impossibility for the declared channel]
No deterministic finite-state controller reading the raw observations of the
witness in Theorem~\ref{thm:rich} decides $\Psi$ in the eventual-output sense.
\end{corollary}

\begin{proof}
The inputs generated by $x_A$ and $x_B$ are the same infinite binary stream.
A deterministic controller therefore has the same state and output sequence on
both states.  That common sequence cannot eventually equal both
$\Psi(x_A)=1$ and $\Psi(x_B)=0$.
\end{proof}

\begin{remark}[Accumulating history versus finite memory]
The history $q_\infty$ is an unbounded record, whereas a finite-state controller
has bounded internal memory.  The history-derived non-repair corollary is the
stronger information restriction; the finite-state result makes the operational
consequence explicit without converting the coarsest decision partition into a
literal memory theorem.
\end{remark}

\section{A repaired two-wall theorem}
\label{sec:wall}

\begin{definition}[Wall-valued target]
A wall-valued orbit target is a map
\[
W:\Xspace\to\{\mathsf{safe},\mathsf{left},\mathsf{right}\}.
\]
Its Boolean safety collapse is
\[
\mathrm{Safe}_W(x)=1\quad\Longleftrightarrow\quad W(x)=\mathsf{safe}.
\]
\end{definition}

\begin{proposition}[Boolean safety does not force wall separation]
A representation sufficient for $\mathrm{Safe}_W$ need not distinguish
$\mathsf{left}$ from $\mathsf{right}$.
\end{proposition}

\begin{proof}
The representation $r=\mathrm{Safe}_W$ decides Boolean safety and places every
failure state in its value-$0$ fiber.
\end{proof}

\begin{proposition}[Wall-valued targets force wall separation]
\label{prop:wall}
If $r:\Xspace\to R$ and $d:R\to\{\mathsf{safe},\mathsf{left},\mathsf{right}\}$
satisfy $W=d\circ r$, then
\[
W(x)\neq W(y)\Longrightarrow r(x)\neq r(y).
\]
\end{proposition}

\begin{proof}
Equal $r$-values would give equal $W$-values after applying $d$.
\end{proof}

\begin{theorem}[Two-wall temporal witness]
There is a history tower for which every $P_n$ descends through $q_n$, while a
wall-valued target $W$ does not descend through $q_\infty$; every representation
sufficient for $W$ separates the left and right failure walls.
\end{theorem}

\begin{proof}
Use
\[
\Xspace=\{A,B_-,B_+\}\times\Bbits^{\mathbb N}\times\mathbb N
\]
with the same $\Endo$, $o$, and $\Pred$ as in Theorem~\ref{thm:rich}.  Define
\[
W(A,s,i)=\mathsf{safe},\quad
W(B_-,s,i)=\mathsf{left},\quad
W(B_+,s,i)=\mathsf{right}.
\]
The three states over $\mathbf 1$ have identical complete observed histories
and different $W$-values, so $W$ fails to descend through $q_\infty$.  The final
claim is Proposition~\ref{prop:wall}.
\end{proof}

\section{A diagnostic is a record-level object}
\label{sec:diagnostic}

\begin{definition}[Diagnostic evidence]
For a pair $(q,\Tail)$, evidence is one of:
\begin{enumerate}[label=\textnormal{(\roman*)},leftmargin=2.6em]
\item a descent certificate $\overline\Tail$ with $\Tail=\overline\Tail\circ q$;
\item a non-descent certificate $(x,y)$ with $q(x)=q(y)$ and $\Tail(x)\neq\Tail(y)$;
\item an incomplete record $\bot$.
\end{enumerate}
\end{definition}

\begin{definition}[Typed diagnostic record]
A diagnostic record is
\[
R=(q,\Tail,e,B),
\]
where $q$ and $\Tail$ have a common domain, $e$ is evidence for $(q,\Tail)$,
and $B$ states the assessment boundary.  Set
\[
\Diag(R)=
\begin{cases}
\DOne,&e\text{ is a valid descent certificate},\\
\DTwo,&e\text{ is a valid non-descent certificate},\\
\DThree,&e=\bot.
\end{cases}
\]
\end{definition}

\begin{proposition}[Diagnostic soundness]
\label{prop:diagnostic}
For a well-typed diagnostic record, $\DOne$ proves $\Dec_q(\Tail)$, $\DTwo$
proves $\neg\Dec_q(\Tail)$, and $\DThree$ asserts only that the record
contains neither certificate.
\end{proposition}

\begin{proof}
The first two claims follow from the definitions and Theorem~\ref{thm:fiber}.
The third follows because $\bot$ is a record status, not a negation of the
existence of a decoder or separating pair.
\end{proof}

\begin{remark}[Excluded readings]
None of $\DOne$, $\DTwo$, or $\DThree$ is a state-level regime, a claim about
all possible representations, or a policy instruction.  In particular,
$\DTwo$ says only that the declared channel is insufficient for the declared
target.
\end{remark}

\section{Binding-sensitive order-two novelty}
\label{sec:novelty}

\begin{definition}[Component library and binding]
Let $\Lib$ be the class of tuples
\[
L=(\Xspace,\Endo,\Pred,\Tail,q^0_\infty,q^1_\infty,x_A,x_B,e_0,e_1)
\]
where $e_0=(x_A,x_B)$ is a valid non-descent certificate for
$(q^0_\infty,\Tail)$ and $e_1$ is a valid descent certificate for
$(q^1_\infty,\Tail)$.  A binding is a tuple
\[
\beta=(q,\Tail,e,B).
\]
An assembled architecture is $A=(L,\beta)$; write $\Arch$ for the class of
such architectures.  The component projection is
\[
\pi_{\mathrm{cmp}}:\Arch\to\Lib,\qquad
\pi_{\mathrm{cmp}}(L,\beta)=L.
\]
\end{definition}

\begin{construction}[Two bindings over one library]
\label{con:bindings}
Use Theorem~\ref{thm:rich}.  Let
\[
q^0_\infty(\theta,s,i)=(s_i,s_{i+1},\ldots),
\qquad
q^1_\infty(\theta,s,i)=(\theta,s_i,s_{i+1},\ldots).
\]
Let $e_0=(x_A,x_B)$ and define
\[
e_1(\theta,t)=
\begin{cases}1,&\theta=A,\\0,&\theta=B.\end{cases}
\]
For the resulting library $L$, define
\[
A_0=(L,(q^0_\infty,\Tail,e_0,B)),\qquad
A_1=(L,(q^1_\infty,\Tail,e_1,B)).
\]
\end{construction}

\begin{theorem}[Binding-sensitive diagnostic non-factorization]
\label{thm:binding}
There is no map $d:\Lib\to\{\DOne,\DTwo,\DThree\}$ such that
\[
\Diag(A)=d\bigl(\pi_{\mathrm{cmp}}(A)\bigr)
\]
for both architectures $A_0$ and $A_1$ of Construction~\ref{con:bindings}.
\end{theorem}

\begin{proof}
The pair $e_0$ yields $\Diag(A_0)=\DTwo$.  The decoder $e_1$ yields
$\Diag(A_1)=\DOne$.  But
$\pi_{\mathrm{cmp}}(A_0)=L=\pi_{\mathrm{cmp}}(A_1)$.  A factorization through
$\pi_{\mathrm{cmp}}$ would therefore imply $\DTwo=\DOne$.
\end{proof}

\begin{definition}[Scoped order-two diagnostic]
Relative to $\pi_{\mathrm{cmp}}$, a diagnostic has representational order two
when it is defined on assembled architectures but does not factor through the
component projection.
\end{definition}

\begin{corollary}[Scoped order-two status]
The diagnostic of Theorem~\ref{thm:binding} has representational order two
relative to $\pi_{\mathrm{cmp}}$.
\end{corollary}

\begin{definition}[Functional novelty]
Let $\Task$ be a typed task, $\Inv$ a protected invariant family, and
$\Base$ a declared baseline class.  The assertion
\[
\FNovel_{\Inv,\Base}(A,\Task)
\]
means that $A$ is well typed and invariant-preserving, realizes $\Task$ on its
stated finite witness class, and that every baseline in $\Base$ fails the task
on that class or exits $\Base$ by retaining a distinction excluded by the
baseline definition.
\end{definition}

\begin{corollary}[Functional novelty relative to componentwise baselines]
Let $\Base_{\mathrm{cmp}}$ consist of outputs that factor through
$\pi_{\mathrm{cmp}}$.  With invariants consisting of common-domain typing,
correct certificate attachment, and the $\DOne/\DTwo/\DThree$ semantics,
$\{A_0,A_1\}$ is a finite non-substitutability witness for the diagnostic task.
\end{corollary}

\begin{proof}
Theorem~\ref{thm:binding} shows that every componentwise baseline returns one
common output on $A_0,A_1$, while the correct outputs differ.
\end{proof}

\begin{remark}[Exact claim]
This is a relative non-preservation result.  It does not claim that familiar
components are individually novel, that unrestricted first-order formalisms
cannot encode the construction, or that historical priority has been established.
\end{remark}

\section{An exact uniform reverse-mathematical principle}
\label{sec:reverse}

\begin{definition}[Gap-tail formulas]
For a coded rational array $a:\mathbb N^2\to\mathbb Q\cap[0,1]$, write
\[
\mathrm{Below}(a,e)\;\Longleftrightarrow\;
\exists k\,\exists M\,\forall m\geq M\;
\bigl(a(e,m)\leq1-2^{-k}\bigr),
\]
\[
\mathrm{AtOne}(a,e)\;\Longleftrightarrow\;
\forall k\,\forall M\,\exists m\geq M\;
\bigl(a(e,m)>1-2^{-k}\bigr),
\]
and
\[
\mathrm{Gap}(a,e)\;\Longleftrightarrow\;
\mathrm{Below}(a,e)\vee\mathrm{AtOne}(a,e).
\]
On a gap-promised row, the two cases express respectively
$\limsup_m a(e,m)<1$ and $\limsup_m a(e,m)=1$.
\end{definition}

\begin{definition}[Uniform gap-tail decision principle]
\[
\UGapTail:\quad
\forall a\left[
\bigl(\forall e\;\mathrm{Gap}(a,e)\bigr)
\Longrightarrow
\exists Y\subseteq\mathbb N\;\forall e\;
\bigl(e\in Y\Longleftrightarrow\mathrm{Below}(a,e)\bigr)
\right].
\]
\end{definition}

\begin{theorem}[Exact strength]
\label{thm:ugap}
Over $\RCA$,
\[
\UGapTail\quad\Longleftrightarrow\quad\ACA.
\]
\end{theorem}

\begin{proof}
Assume $\ACA$.  The formula $\mathrm{Below}(a,e)$ is arithmetical with
parameter $a$, so arithmetical comprehension produces
$Y=\{e:\mathrm{Below}(a,e)\}$.

Conversely, assume $\UGapTail$.  It suffices over $\RCA$ to obtain the range
of every injection.  Let $f:\mathbb N\to\mathbb N$ be injective and define
\[
a_f(y,m)=
\begin{cases}
1-\dfrac{1}{m+1},&\exists x\leq m\;f(x)=y,\\[1ex]
0,&\text{otherwise.}
\end{cases}
\]
The array is uniformly recursive in $f$.  If $y\in\ran(f)$, the row is
eventually $1-1/(m+1)$ and satisfies $\mathrm{AtOne}(a_f,y)$.  If
$y\notin\ran(f)$, it is identically $0$ and satisfies
$\mathrm{Below}(a_f,y)$.  Thus all rows are gap-promised and
\[
\mathrm{Below}(a_f,y)\Longleftrightarrow y\notin\ran(f).
\]
Applying $\UGapTail$ gives a set for the complement of $\ran(f)$; recursive
comprehension gives its complement.  The range-existence principle for
injections is equivalent over $\RCA$ to $\ACA$.
\end{proof}

\begin{remark}[No inflation]
Theorem~\ref{thm:ugap} concerns the stated uniform family only.  It does not
make a claim about the single temporal witness, every tail predicate, or any
application of the diagnostic.
\end{remark}


\section{Universal M-space proof architecture}
\label{sec:universal-mspace}

The preceding results decide when a declared target survives a declared
representation.  This section applies that same factorization discipline to
mathematical proofs themselves.  It formalizes four architectural claims:
universal representation on the maximal finitely communicable proof domain;
maximal compression relative to a declared audit family; three-dimensional
visual realization of the resulting audit quotient; and necessity, up to
audit-equivalent representation, in novelty-separated derivational domains.
The detailed constructions and full proofs are supplied in the supplement.

\begin{definition}[Finitarily checkable calculus and communicable proof witness]
\label{def:communicable-v9}
A finitarily checkable calculus is a finite code
\[
C=(\Sigma_C,\mathsf{Form}_C,\mathsf{Rule}_C,\mathsf{End}_C,\mathsf{Check}_C)
\]
in which formulas are decidable finite strings, inference records are finite
directed acyclic hypergraphs with finite premise lists, and
$\mathsf{Check}_C$ is a terminating local verifier.  A communicable purely
mathematical proof witness is a pair $p=(C,D)$ in which $D$ is an accepted
finite derivation record.  A recursively described proof is included when its
generator, boundary conditions, and local verifier have a finite accepted code.
Write $\Proofs$ for this class.
\end{definition}

The finite-communication restriction is exact rather than discretionary: no
finite artifact language can injectively represent an uncountable class of
distinct non-finitely-specifiable objects.  Thus $\Proofs$ is the largest
domain on which a theorem of universal finite communicable representation can
be true.

\begin{definition}[M-space proof tower]
\label{def:mtower-v9}
For $n\in\mathbb N$, let $\mathcal M_n$ be the class of finite tuples
\[
m=(G,L_0,L_1,\ldots,L_n),
\]
where $G$ is a finite typed directed proof-incidence graph and each $L_i$ is a
finite typed record layer attached to vertices, edges, paths, or finite
subcomplexes.  The base layers store respectively exact source syntax,
normalized type information, local inference and binding certificates, and
audit/provenance/display obligations; higher layers store finite certificates
about lower-layer certificates.  The lifts append an empty layer and the
projections delete the last layer.  Set
\[
\MInf=\bigcup_{n\in\mathbb N}\mathcal M_n
\]
under these lifts.
\end{definition}

\begin{theorem}[Universal M-space proof representation]
\label{thm:universal-v9}
There are maps
\[
\MEnc:\Proofs\longrightarrow\MInf,
\qquad
\MDec:\operatorname{im}(\MEnc)\longrightarrow\Proofs
\]
such that for every $p\in\Proofs$ there is a finite $N(p)$ with
$\MEnc(p)\in\mathcal M_{N(p)}$ and
\[
\MDec(\MEnc(p))=p.
\]
The typed incidence graph, ordered premise bindings, local checker records,
certificate references, provenance records, and declared audit obligations are
recoverable from $\MEnc(p)$.
\end{theorem}

\begin{proof}
Store the exact finite code $(C,D)$ in $L_0$, construct the finite typed
incidence graph of $D$, and attach the finite normalization, certificate,
provenance, and audit records in successive finite layers.  The decoder reads
the exact $L_0$ serialization and verifies agreement with the derived layers.
All construction steps terminate on finite data, so the resulting object lies
at a finite level and decodes exactly to $p$.
\end{proof}

\begin{definition}[Audit family and audit quotient]
\label{def:auditquot-v9}
An audit family $\mathcal A$ on $\Proofs$ is a set of task maps
$a:\Proofs\to Y_a$.  Define
\[
p\AuditEq p'
\quad\Longleftrightarrow\quad
\forall a\in\mathcal A\,[a(p)=a(p')],
\qquad
\AuditQuot(p)=[p]_{\AuditEq}.
\]
A representation $r:\Proofs\to R$ is audit-faithful when every
$a\in\mathcal A$ factors through $r$.
\end{definition}

\begin{theorem}[Audit-optimal compression]
\label{thm:auditquot-v9}
The audit quotient $\AuditQuot$ is audit-faithful.  For every
audit-faithful representation $r:\Proofs\to R$, there is a unique map
\[
h_r:\operatorname{im}(r)\to\Proofs/{\AuditEq}
\]
such that
\[
\AuditQuot=h_r\circ r.
\]
Hence $\AuditQuot$ is the coarsest representation preserving every declared
audit distinction.
\end{theorem}

\begin{proof}
Each task $a$ is well defined on an $\AuditEq$-class, so it factors through
$\AuditQuot$.  If $r(p)=r(p')$, audit faithfulness implies
$a(p)=a(p')$ for every $a$, hence $p\AuditEq p'$.  Therefore
$h_r(r(p))=[p]_{\AuditEq}$ is well defined and is unique on the image.
\end{proof}

The word ``maximal'' consequently has a precise invariant meaning: the quotient
removes every distinction irrelevant to the declared audit and no distinction
needed by it.  It does not claim an impossible universal dominance in raw
description length across all lossless codes.

\begin{lemma}[Universal edge-separated three-dimensional realization]
\label{lem:3d-v9}
Every finite proof-dependency graph has an edge-separated labelled realization
in $\Rthree$.
\end{lemma}

\begin{proof}
Assign distinct parameters $t_1,\ldots,t_n$ to the vertices and place vertex
$i$ at $(t_i,t_i^2,t_i^3)$.  No four distinct points on this moment curve are
coplanar.  If two nonincident straight edges intersected, their four endpoints
would be coplanar, a contradiction.  Finite labels preserve vertex and edge
types, premise order, and certificate references.
\end{proof}

\begin{theorem}[Three-dimensional universal minimality]
\label{thm:3dmin-v9}
Three is the least Euclidean dimension in which every finite
proof-dependency graph admits an edge-separated visual realization.
\end{theorem}

\begin{proof}
Lemma~\ref{lem:3d-v9} proves sufficiency.  For necessity, a finite proof ledger
can contain a subdivision of $K_5$: take five certified axiom occurrences and,
for each pair, a conjunction-introduction inference joining the two occurrences.
An edge-separated realization in $\mathbb R^2$ would make this subdivision
planar, hence make $K_5$ planar after contraction, contradicting
$|E(K_5)|=10>3\cdot5-6$.  Dimensions zero and one fail a fortiori.
\end{proof}

For graph-shaped proof-dependency audit, the quotient has a single labelled
edge-separated $\Rthree$ realization.  Its use as a single-view audit display is
conditional on the full labelled incidence graph satisfying the declared
finite-interface capability bound.  For audit content carried by intrinsic
geometry rather than dependency incidence, the supplement proves the
corresponding conditional result for a complete aspect--scale decomposition
equipped with a finite admissible projection plan: each proof activates only a
finite indexed record of coherent three-dimensional views, supplied with local
and global decoders that recover the declared audit quotient.  Under the
explicit finite-legend and trusted-local-kernel capability model, that record
supports a finite sound-and-complete structural audit without theorem-specific
subject-matter expertise.  This is a conditional interface theorem, not an
assertion about unmodeled cognition.

\begin{definition}[Novelty-separated derivational problem]
\label{def:novel-v9}
A derivational problem is a triple $\mathfrak D=(X,T,\mathcal R)$, where
$X\subseteq\Proofs$ is a declared class of proof or proof-search states,
$T:X\to Y$ is a required derivational or audit target, and $\mathcal R$ is a
family of available representation channels.  It is novelty-separated relative
to $\mathcal R$ when
\[
\forall r\in\mathcal R\;\exists x,y\in X
\quad r(x)=r(y)\ \text{and}\ T(x)\neq T(y).
\]
\end{definition}

\begin{theorem}[Representation necessity up to audit equivalence]
\label{thm:necessity-v9}
No channel in a novelty-separated family supports a reliable decision of $T$ on
$X$.  If $T$ includes a full audit family $\mathcal A$, every reliable
replacement must retain a refinement through which the restricted quotient
$\AuditQuot|_X$ factors.  When the required audit is universal,
edge-separated, and visually traceable by the declared general-background
auditor, the replacement must additionally admit an audit-equivalent
three-dimensional factor: one view for graph-shaped content and a coherent
family of views otherwise.
\end{theorem}

\begin{proof}
The first statement is the fiber criterion: equal representation values with
unequal target values preclude factorization of $T$.  The second is
Theorem~\ref{thm:auditquot-v9} applied to $X$.  The last statement follows from
Theorem~\ref{thm:3dmin-v9} and from the finite admissible projection-plan
construction given in the supplement.  A singleton visual factor is available
for graph-shaped content only when the full selected incidence display satisfies
the declared interface bound.
\end{proof}


\subsection{Collective proof-audit capacity and empirical bridge}
\label{sec:capacity-v10}

The preceding necessity theorem is pointwise: a representation that merges
target-distinct proof states cannot decide the target.  A collective mathematical
workflow introduces a second question: whether the rate at which source-visible
audit obligations arrive can be resolved with the audit capacity available to a
community.  The following ledger formalizes that question without identifying a
repository, a field, or a bibliographic network with mathematics as a whole.

\begin{definition}[Time-indexed proof-audit ledger]
\label{def:ledger-v10}
For discrete windows $t\in\mathbb N$, a proof-audit ledger is
\[
  \Ledger_t=(V_t,E_t,O_t,R_t,\tau_t).
\]
Here $V_t$ is a finite set of versioned proof-relevant artifacts, $E_t$ is a
finite typed relation set (including, when supported by the source, dependency,
review, challenge, repair, supersession, and resolution), $O_t\subseteq V_t$ is
the unresolved audit-item set, $R_t$ is the set of recorded resolutions, and
$\tau_t$ records source version and provenance.  This is a ledger of
source-visible artifacts, not a claim of ontological completeness for
mathematics.
\end{definition}

\begin{definition}[Audit-work accounting]
\label{def:work-v10}
Fix an audit family $\mathcal A$ and an $\mathcal A$-faithful representation
$\rho$.  Let $A_t$ be the finite set of audit obligations opened or reopened in
window $t$, let $c_\rho(o)\geq 0$ be the work charged to resolve $o$ through
$\rho$, and let $\AuditCap_t\geq0$ be the work capacity available in that window.
Define the incoming audit demand and debt recurrence by
\[
 \Demand_\rho(t)=\sum_{o\in A_t}c_\rho(o),\qquad
 \Debt_{\rho,t+1}=\max\{0,\Debt_{\rho,t}+\Demand_\rho(t)-\AuditCap_t\}.
\]
The units of $c_\rho$ may be time, verified traversal steps, or a declared
weighted audit count; they must be fixed before any empirical use.  This is a
deterministic reflected workload recurrence; no stochastic queueing premise is
used by the theorem below.
\end{definition}

\begin{assumption}[Audit-cost monotonicity]
\label{ass:cost-v10}
For audit-faithful representations $\rho_1$ and $\rho_2$, if
$\rho_1=u\circ\rho_2$ for some map $u$ on $\operatorname{im}(\rho_2)$, then
$c_{\rho_1}(o)\leq c_{\rho_2}(o)$ for every obligation $o$ in the declared
domain.  Thus retaining audit-irrelevant distinctions cannot reduce the declared
structural audit work.
\end{assumption}

\begin{theorem}[Audit-debt divergence]
\label{thm:debt-v10}
If there are $t_0\in\mathbb N$ and $\varepsilon>0$ such that
\[
  \Demand_\rho(t)-\AuditCap_t\geq\varepsilon
  \quad\text{for every }t\geq t_0,
\]
then
\[
  \Debt_{\rho,t}\geq\Debt_{\rho,t_0}+(t-t_0)\varepsilon
  \quad\text{for every }t\geq t_0,
\]
and hence $\Debt_{\rho,t}\to\infty$.
\end{theorem}

\begin{proof}
The recurrence in \Cref{def:work-v10} gives
$\Debt_{\rho,t+1}\geq\Debt_{\rho,t}+\varepsilon$ for $t\geq t_0$.
Induction proves the displayed inequality.
\end{proof}

\begin{theorem}[Audit-quotient capacity dominance]
\label{thm:capacity-v10}
Under \Cref{ass:cost-v10}, for every $\mathcal A$-faithful representation
$r$ and every time window,
\[
  \Demand_{\AuditQuot}(t)\leq\Demand_r(t).
\]
If the two ledgers begin with equal debt and share the same capacity schedule,
then
\[
  \Debt_{\AuditQuot,t}\leq\Debt_{r,t}
  \qquad(t\in\mathbb N).
\]
Suppose that $r$ has a persistent positive demand-capacity gap while
\[
  \Demand_{\AuditQuot}(t)-\AuditCap_t\leq0
  \qquad(t\geq t_0)
\]
for some $t_0$.  Then replacement of $r$ by the audit quotient is a sufficient
representation-level repair: quotient debt is nonincreasing after $t_0$, while
the debt of $r$ diverges.

More generally, avoiding the divergence conclusion requires changing at least
one stated antecedent of \Cref{thm:debt-v10}: the representation-specific
demand, the capacity schedule, the incoming obligation process, the declared
audit domain or task, or the persistence condition.  This does not assert that
every successful repair must literally use the quotient realization.
\end{theorem}

\begin{proof}
By \Cref{thm:auditquot-v9}, $\AuditQuot=h_r\circ r$ on
$\operatorname{im}(r)$.  Audit-cost monotonicity therefore gives
$c_{\AuditQuot}(o)\leq c_r(o)$ for each incoming obligation, and summation gives
the demand inequality.  The debt inequality follows by induction from the
monotonicity of $x\mapsto\max\{0,x\}$.

Under the nonpositive quotient excess condition,
\[
  \Debt_{\AuditQuot,t+1}\leq\Debt_{\AuditQuot,t}
  \qquad(t\geq t_0).
\]
The persistent excess of $r$ yields divergence by \Cref{thm:debt-v10}.  The
final statement follows because that divergence theorem depends only on its
displayed recurrence and antecedent conditions; another repair may avoid
divergence by changing one of them without literally adopting the quotient.
\end{proof}

\begin{remark}[Formal criterion versus empirical bridge]
\label{rem:protocol-v10}
\Cref{thm:debt-v10,thm:capacity-v10} is a theorem about a declared ledger and
work model.  It does not assert that an informal mathematical field presently
has a growing debt, that any particular graph statistic measures truth, or that
a three-dimensional interface is empirically fastest for every user.  Those
questions are bridge-required.  A companion preregistered protocol specifies
source-visible formal corpora (for example, \texttt{mathlib} and the Archive of
Formal Proofs), locked-test forecasting of future ledger debt, and a randomized
comparison of text, 2D, quotient-linked 2D, and quotient-linked 3D structural
audit conditions \cite{mathlib2020,vanDoorn2020,blanchette2015,huch2024,nosek2018}.
Bibliographic graphs are secondary communication evidence rather than proof
validity evidence; algebraic connectivity is a descriptive graph statistic, not
an epistemic truth measure \cite{fiedler1973,chung1997}.
\end{remark}

\begin{theorem}[Assembled universal M-space proof architecture]
\label{thm:fullarch-v9}
For the class $\Proofs$ of communicable purely mathematical proof witnesses:
\begin{enumerate}[label=\textnormal{(\roman*)},leftmargin=2.6em]
\item every proof has an exact finite M-space representation and decoder;
\item the audit quotient is maximally compressive in the task-relative quotient
sense;
\item given a complete aspect--scale decomposition and a finite admissible
projection plan, the quotient has an audit-preserving finite indexed record of
edge-separated three-dimensional views; for graph-shaped proof-dependency audit,
this record is a singleton exactly when the full labelled incidence display
satisfies the declared interface bound;
\item three dimensions are universally minimal for edge-separated proof-graph
tracing;
\item under the declared structural-auditor model, a plan-realized structural
audit family, a terminating local kernel, and a finite admissible projection
plan, the displayed record is sound and complete for the declared structural
audit;
\item in novelty-separated derivational domains, reliable derivation requires an
audit-equivalent refinement of the quotient and, under the visual-audit
requirement, an audit-equivalent three-dimensional factor realized by the
declared finite plan; and
\item under the declared audit-work accounting and audit-cost monotonicity,
the audit quotient weakly minimizes incoming audit demand among faithful
representations.  A persistent representation-induced demand-capacity gap
produces unbounded debt; if the quotient has nonpositive excess, quotient
replacement is a sufficient representation-level repair, without excluding
other repairs that change a stated premise of the recurrence.
\end{enumerate}
\end{theorem}

\begin{proof}
Items (i)--(iii) are Theorems~\ref{thm:universal-v9} and
\ref{thm:auditquot-v9} together with the finite admissible projection-plan
construction in the supplement; item (iv) is Theorem~\ref{thm:3dmin-v9};
item (v) is the supplement's corrected conditional structural-audit theorem;
item (vi) is Theorem~\ref{thm:necessity-v9}; and item (vii) is
Theorem~\ref{thm:capacity-v10}.
\end{proof}

\begin{remark}[Scope of the architectural theorem]
\label{rem:scope-v9}
The theorem is about finite communicable proof witnesses, task-relative audit
compression, and structural derivation/audit.  It does not assert a universal
algorithm for discovering all mathematics, semantic soundness of an arbitrary
calculus without a soundness certificate, a literal-format monopoly, or an
unmodeled empirical claim that every person can understand every mathematical
concept merely by viewing a display.  It also does not claim that a source-visible
ledger measures all mathematical activity.  The invariant conclusion is stronger
and more exact: any reliable workflow for the stated task must preserve the audit
quotient; universal edge-separated visual tracing requires a
three-dimensional-equivalent factor realized under its declared finite audit
plan; and retaining audit-irrelevant distinctions cannot improve the formal
demand margin under the stated cost model.  Quotient replacement is sufficient
when it removes the persistent excess, but the theorem does not impose a literal
monopoly on all successful repairs.
\end{remark}

\section{Applications, auditability, and bridge discipline}
\label{sec:bridge}

\begin{definition}[Application bridge record]
An application bridge record is
\[
\mathsf{Bridge}_{\mathrm{app}}=
\bigl(X_{\mathrm{app}},\iota,T_{\mathrm{app}},o_{\mathrm{app}},
\Psi_{\mathrm{app}},\mathcal R,\mathcal E,\mathsf{Status}\bigr),
\]
where $\iota:X_{\mathrm{app}}\to X$ is a declared encoding and the record
separately states whether
\[
\iota\circ T_{\mathrm{app}}=T\circ\iota,\qquad
o_{\mathrm{app}}=o\circ\iota,\qquad
\Psi_{\mathrm{app}}=\Psi\circ\iota
\]
hold on the intended domain.  The lists $\mathcal R$ and $\mathcal E$ record
unresolved premises and defeater forms.  The status set is
\[
\{\mathsf{THM},\mathsf{CND},\mathsf{BR},\mathsf{HEUR},\mathsf{OUT},\mathsf{DEF}\}.
\]
\end{definition}

The M-space proof architecture of Section~\ref{sec:universal-mspace} is an
additional formal core, independent of Theorems~\ref{thm:fiber}--\ref{thm:ugap}.
It proves exact representation for communicable proof witnesses, quotient-optimal
audit compression, and the minimal universal dimension for edge-separated
proof-dependency visualization.  It proves necessity only up to
audit-equivalent representation and only under its declared novelty-separation,
finite-plan, and visual-audit requirements.  Its ledger-capacity extension is
likewise conditional on the declared work model; the companion corpus and
interface protocol is required to test its empirical premises.  None of these results makes
an application bridge true, establishes the semantic soundness of an arbitrary
model without its own soundness certificate, or entails a physical,
institutional, or policy conclusion.

\section{Conclusion}

For a declared observation map, an orbit target is decidable exactly when it is
fiber-constant.  The temporal witness shows that every declared finite monitor
may be decidable at its matching stage while a distinct long-run target remains
undecidable from full observed history.  The correct consequence is a
record-level diagnostic of representational adequacy, not a state-level regime or
an action policy.

The binding-sensitive theorem makes the paper's functional claim visible: a
component list can preserve all local ingredients yet lose the channel--target--
certificate relation needed for the diagnostic.  The M-space theorem package
extends this point to proof architecture: every communicable proof has an exact
typed M-space encoding; the audit quotient is the coarsest task-faithful
compression; and the quotient has a universally dimension-minimal
edge-separated three-dimensional realization for proof-dependency audit.  In
novelty-separated derivational domains, a reliable workflow must preserve this
audit content up to representation equivalence.  The proof-audit ledger adds a
collective constraint: under a stated cost model, an audit quotient weakly
minimizes workload, and a persistent excess of workload over capacity produces
unbounded unresolved audit debt.  When the quotient itself has nonpositive
excess, quotient replacement is a sufficient representation-level repair;
other repairs are not excluded if they change a stated recurrence premise.
Whether formal repositories exhibit that pattern, and whether quotient-linked
2D or 3D displays improve actual audit performance, is deliberately allocated
to a preregistered empirical bridge.
Separately, $\UGapTail\leftrightarrow\ACA$ gives an exact proof-theoretic theorem
for a uniform family.  Every external application or physical interpretation
remains bridge-required.

\begin{thebibliography}{99}
\bibitem{AlpernSchneider1985}
B. Alpern and F. B. Schneider.
\newblock Defining liveness.
\newblock \emph{Information Processing Letters}, 21(4):181--185, 1985.

\bibitem{Pnueli1977}
A. Pnueli.
\newblock The temporal logic of programs.
\newblock In \emph{Proceedings of the 18th Annual Symposium on Foundations of
Computer Science}, pages 46--57. IEEE, 1977.

\bibitem{Simpson2009}
S. G. Simpson.
\newblock \emph{Subsystems of Second Order Arithmetic}, second edition.
\newblock Cambridge University Press, 2009.

\bibitem{blanchette2015}
J.~C. Blanchette, M.~Haslbeck, D.~Matichuk, and T.~Nipkow.
\newblock Mining the Archive of Formal Proofs.
\newblock In \emph{Intelligent Computer Mathematics}, pages 3--17. Springer, 2015.
\newblock doi: \href{https://doi.org/10.1007/978-3-319-20615-8_1}{10.1007/978-3-319-20615-8\_1}.

\bibitem{chung1997}
F.~R.~K. Chung.
\newblock \emph{Spectral Graph Theory}.
\newblock American Mathematical Society, 1997.

\bibitem{fiedler1973}
M.~Fiedler.
\newblock Algebraic connectivity of graphs.
\newblock \emph{Czechoslovak Mathematical Journal}, 23(2):298--305, 1973.
\newblock doi: \href{https://doi.org/10.21136/CMJ.1973.101168}{10.21136/CMJ.1973.101168}.

\bibitem{huch2024}
F.~Huch and M.~Wenzel.
\newblock Distributed parallel build for the Isabelle Archive of Formal Proofs.
\newblock In \emph{15th International Conference on Interactive Theorem Proving (ITP 2024)},
LIPIcs 309, Article 22, 2024.
\newblock doi: \href{https://doi.org/10.4230/LIPIcs.ITP.2024.22}{10.4230/LIPIcs.ITP.2024.22}.

\bibitem{mathlib2020}
The mathlib Community.
\newblock The Lean mathematical library.
\newblock In \emph{Proceedings of the 9th ACM SIGPLAN International Conference on Certified Programs and Proofs},
pages 367--381, 2020.
\newblock doi: \href{https://doi.org/10.1145/3372885.3373824}{10.1145/3372885.3373824}.

\bibitem{nosek2018}
B.~A. Nosek, C.~R. Ebersole, A.~C. DeHaven, and T.~Mellor.
\newblock The preregistration revolution.
\newblock \emph{Proceedings of the National Academy of Sciences}, 115(11):2600--2606, 2018.
\newblock doi: \href{https://doi.org/10.1073/pnas.1708274114}{10.1073/pnas.1708274114}.

\bibitem{vanDoorn2020}
F.~van Doorn, G.~Ebner, and R.~Y. Lewis.
\newblock Maintaining a library of formal mathematics.
\newblock arXiv:2004.03673, 2020.
\end{thebibliography}

\end{document}
